\chapter{Database Design}
\label{ch:database}

\section{Database Overview}
The application uses a MySQL relational database with two tables:
\texttt{users} and \texttt{game\_results}. The schema is intentionally
small, reflecting the project's single-role, two-entity data model.

\section{Tables}

\begin{table}[h]
\centering
\caption{\texttt{users} table.}
\label{tab:users-table}
\begin{tabular}{@{}p{2.6cm}p{3.4cm}p{7.5cm}@{}}
\toprule
\textbf{Column} & \textbf{Type} & \textbf{Description} \\
\midrule
id & INT, PK, AUTO\_INCREMENT & Unique user identifier \\
username & VARCHAR, UNIQUE & Display name, 3--50 characters \\
email & VARCHAR, UNIQUE & Login email address \\
password\_hash & VARCHAR & Bcrypt password hash via \texttt{password\_hash()} \\
created\_at & DATETIME & Account creation timestamp \\
\bottomrule
\end{tabular}
\end{table}

\begin{table}[h]
\centering
\caption{\texttt{game\_results} table.}
\label{tab:results-table}
\begin{tabular}{@{}p{3cm}p{3.4cm}p{7cm}@{}}
\toprule
\textbf{Column} & \textbf{Type} & \textbf{Description} \\
\midrule
id & INT, PK, AUTO\_INCREMENT & Unique game record identifier \\
user\_id & INT, FK $\to$ users(id) & Owning player \\
difficulty & ENUM & \texttt{beginner} or \texttt{advanced} \\
score & INT & Server-recomputed final score \\
time\_taken & INT & Seconds survived (0--120) \\
cells\_revealed & INT & Number of safe cells revealed \\
hints\_used & INT & Number of hints used (0--3) \\
created\_at & DATETIME & Result submission timestamp \\
\bottomrule
\end{tabular}
\end{table}

\section{Relationships}
The relationship between \texttt{users} and \texttt{game\_results} is a
standard one-to-many relationship: one user can accumulate many game
records over time, so \texttt{user\_id} is stored as a foreign key on the
\texttt{game\_results} side.

\begin{figure}[h]
\centering
\begin{tikzpicture}[node distance=3cm]
    \tikzstyle{entity} = [rectangle, draw=black, thick, fill=blue!8,
        text width=3.2cm, align=center, minimum height=1cm]
    \node (users) [entity] {users};
    \node (results) [entity, right=of users] {game\_results};
    \draw [-{Latex[length=2mm]}, thick] (users) -- node[above] {1 -- *} (results);
\end{tikzpicture}
\caption{Entity relationship between \texttt{users} and \texttt{game\_results}.}
\label{fig:er-diagram}
\end{figure}

\section{Important Constraints}
\begin{itemize}[nosep]
    \item Foreign key: \texttt{game\_results.user\_id} references
    \texttt{users.id}.
    \item Unique constraints on \texttt{users.username} and
    \texttt{users.email} prevent duplicate accounts.
    \item ENUM constraint on \texttt{game\_results.difficulty} restricts the
    column to exactly \texttt{beginner} or \texttt{advanced}.
    \item An index is defined on \texttt{game\_results.user\_id} (and on
    \texttt{score} for leaderboard ordering) to support efficient history
    and ranking queries.
\end{itemize}

\section{Schema Evolution}
An early version of the schema included a \texttt{result} column recording
a win/loss outcome, consistent with a classic ``clear the board'' win
condition. When the game was redesigned as a timed survival mode
(Section~\textit{Survival Mode and the Absence of a Win Condition},
Chapter~\ref{ch:game-design}), this column was no longer meaningful and was
removed via \texttt{migration-remove-win-condition.sql}, since a completed
game is now fully described by its difficulty, time survived, cells
revealed, and hints used.
